\documentclass[11pt]{article}
\usepackage[margin=1in]{geometry}
\usepackage{amsmath,amssymb,amsthm,hyperref}
\newtheorem{definition}{Definition}
\newtheorem{theorem}{Theorem}
\newtheorem{lemma}{Lemma}
\newtheorem{corollary}{Corollary}
\newtheorem{obligation}{Open obligation}
\newcommand{\SAT}{\mathsf{SAT}}
\newcommand{\PP}{\mathsf P}
\newcommand{\NP}{\mathsf{NP}}
\title{LOWER\_BOUND through LB-0015: Universal Quantifiers, Clean-Core Guards, and the Oracle Barrier}
\author{P versus NP project}
\date{11 October 2026}
\begin{document}
\maketitle
\paragraph{Status.}
This manuscript proves elementary model-audit lemmas and a matched-oracle
extension of the prefix-search immunity argument in Hirahara
\cite{hirahara2020}. It closes a barrier obligation deferred in project
PVG-4-0001, not the existence obligation for PFCERT. No separation of
$\PP$ and $\NP$ is claimed. These are human-readable formal proofs,
not proof-assistant verification.

\section{What history distinguishability actually proves}
\begin{definition}[Summary-only continuation]
Fix a set $H$ of histories, a common query set $Q$, and a response function
$f:H\times Q\to\{0,1\}$. A summary-only implementation consists of
$s:H\to\{0,1\}^{b}$ and one function
$E:\{0,1\}^{b}\times Q\to\{0,1\}$ satisfying
$E(s(h),q)=f(h,q)$. No history-dependent source, advice, oracle, or
unrecorded state is available to $E$.
\end{definition}
\begin{lemma}[Summary bound]
If $H_0\subseteq H$ has $M$ pairwise distinct response rows
$f(h,\cdot)$, then every summary-only implementation has
$b\ge\lceil\log_2M\rceil$.
\end{lemma}
\begin{proof}
If $s(h)=s(h')$, the common evaluator has the same answer for every
query, so the two rows are equal. Thus $s$ is injective on $H_0$.
\end{proof}
This proves neither $M$ time steps nor a bound for a model allowing
source rereading. For example, let $H=\{0,1\}^{k}$ and
$f(x,i)=x_i$. All $2^k$ rows differ, but if $x$ remains on the read-only
random-access input tape, each query requires just its indexed input
read and ordinary address handling. There need be no retained $k$-bit
copy of $x$ in work space.

\begin{lemma}[Endpoint bypass, recovered]
Assume $A$ is an exact polynomial-time SAT decider. If the sole required
output on a formula $F$ is $\SAT(F)$ and the ordinary model permits
access to $F$ and execution of $A$, then there is a correct implementation
that ignores every wrapper-created history and returns $A(F)$.
\end{lemma}
\begin{proof}
The stated implementation has the required output and the assumed
polynomial time bound by the definition of $A$.
\end{proof}
Consequently, an auxiliary witness being unique, or a false history
destroying satisfying extensions of a wrapper CNF, does not prove that
an arbitrary decision algorithm constructs that witness or history.

\section{The current all-input gap}
Let $\mathcal A$ be the set of simultaneous SAT time-space exponent
pairs for one uniform all-length deterministic machine on the declared
encoding and random-access model. The project has proved a conditional
recursive composition rule
\[
 C=\max\left\{1,r,\frac{v-h}{1-h}\right\},\qquad
 P=\max\{p,h\},
\]
for a uniformly hereditary, fully charged all-input phase with local
time $N^{r+o(1)}$, aggregate executed child volume $N^{v+o(1)}$,
maximum child length $N^{h+o(1)}$, $h<1$, and local state
$N^{p+o(1)}$. Repeated children are counted with multiplicity.
The proof and full contract are in FORMAL-SPINE1 and
RECOVERY-STRENGTHEN7, identified in the companion audit report.

\begin{lemma}[Existential phase equivalence, recovered]
Suppose a phase contract permits zero children and local direct decision
at resource bounds $(C,P)$, and recursive composition converts every
such phase to a SAT algorithm with the same bounds. Then existence of
an all-input phase is equivalent to $(C,P)\in\mathcal A$.
\end{lemma}
\begin{proof}
The forward implication is the assumed composition theorem.
For the reverse implication, execute the algorithm locally and issue no
children. The child restrictions hold vacuously.
\end{proof}
A positive child count alone does not repair this: one can issue and
ignore a constant SAT instance.

\begin{obligation}[Source-visible all-input construction]
Prove, under the assumption that SAT has a polynomial-time decider, the
existence of a uniform all-length implementation satisfying a favorable
fully charged hereditary phase contract on every root and actual
unresolved descendant. A sufficient stronger target is an explicit
compiler from a clocked SAT decider to a PRE-A construction that builds
strictly smaller, semantically essential children before any parent-scale
semantic call. Acquisition failures, replay, input access, current child
buffers, and all ancestor state must be charged in the same machine.
\end{obligation}
This obligation is separation-strength for contracts whose endpoint is
already forbidden. The following verified endpoint is a checkpoint,
not a solution to the obligation.

Fortnow--Lipton--van Melkebeek--Viglas \cite{flmv2004}, Corollary 5.8
and the local SAT reduction, exclude $(t,s)$ with strict margins
\[
 t<c,\qquad
 s<d(c):=\frac12\left(\frac{c+2}{c^2}-c\right).
\]
For $c=1201/1000$,
\[
 d(c)=\frac12+\frac{26275399}{2884802000},
\]
so $(6/5,1/2)$ is forbidden. Matching representation, local reduction,
and simultaneous resource bounds are essential. This is not an
unrestricted polynomial-time SAT lower bound.

\section{Matched oracle and machine conventions}
Fix an arbitrary oracle $O\subseteq\{0,1\}^{*}$ and a universal binary-program
oracle machine $U^O$. Every computation and verifier in this section uses
the same $O$. A program has at most one final output, produced when it
halts. Query writing and ordinary tape work are counted; the oracle
answer step has the usual unit cost.

We use a plain universal machine and require efficient universality:
for each fixed plain-description oracle machine $M$
there is a compiler prefix of constant length independent of its input,
and simulation by $U^O$ has polynomial overhead in input length and
running time, uniformly in $O$. No constant-overhead claim for simulating
arbitrary plain-description machines by a prefix-free universal machine
is made. A self-delimiting encoding of
$n$ of length at most $2\lceil\log_2n\rceil+O(1)$ is used.

Let $t:\mathbb N\to\mathbb N$ eventually dominate every polynomial.
Define
\[
 K_U^{t,O}(x)=\min\{|p|:U^O(p)\text{ halts with final output }x
                   \text{ within }t(|x|)\text{ steps}\},
\]
with minimum $+\infty$ when no such program exists.

\begin{theorem}[Relativized immunity with logarithmic deficiency]
Let $b(n)$ be a nonnegative integer with
$b(n)\le c\log_2(n+2)$ for some constant $c$. Then
\[
 R_b^{t,O}=\{x:K_U^{t,O}(x)\ge |x|-b(|x|)\}
\]
is $\PP^O$-immune.
\end{theorem}
\begin{proof}
Counting programs of length strictly less than $n-b(n)$ shows that,
for every sufficiently large $n$, at least one $n$-bit string lies in
$R_b^{t,O}$. Thus the set is infinite.

Suppose an infinite $L\in\PP^O$ were contained in $R_b^{t,O}$.
Fix its deterministic polynomial-time decider $M^O$.
Choose a positive integer constant $a>c+4$ and set
$m(n)=\lceil a\log_2(n+2)\rceil$.
All sufficiently large $n$ have $m(n)<n$; ignore the finitely many others.

For any sufficiently large length $n$ at which $L$ is nonempty, choose
one $y\in L\cap\{0,1\}^n$ and let $u$ be its first $n-m(n)$ bits.
A fixed oracle program, given encoded $n$ and $u$, enumerates all
$m(n)$-bit suffixes $v$, tests $M^O(uv)$, and prints the first accepted
$uv$. Such a suffix exists. Its running time, including universal
simulation, is polynomial in $n$: there are
$2^{m(n)}=n^{O(1)}$ tests of a fixed polynomial-time machine. This time
is eventually at most $t(n)$.

The program description has length at most
\[
 n-m(n)+2\lceil\log_2 n\rceil+O(1)
 <n-b(n)
\]
for all sufficiently large $n$. The selected output $x=uv$ belongs
to $L$, but has $K_U^{t,O}(x)<n-b(n)$, a contradiction.
\end{proof}

The proof is a direct matched-oracle adaptation and logarithmic-deficiency
extension of the published prefix-search argument, not a new
nonrelativizing mechanism.

\section{A generalized bounded-list target}
Let $L(n)\ge1$ be an integer, polynomially bounded and polynomial-time
computable. Define
\[
 \ell(n)=\lceil\log_2 L(n)\rceil,\qquad
 k(n)=n-\ell(n)-1.
\]
Only sufficiently large $n$, for which $k(n)\ge0$, are relevant.

\begin{definition}[Polynomially witnessed endpoint list cover]
One fixed deterministic polynomial-time oracle verifier
$V^O(1^n,p,x,\pi)$ and one fixed polynomial $q$ satisfy:
\begin{description}
\item[] \emph{Witness convention.} Every accepted witness is restricted
to $|\pi|\le q(n)$; outputs have $|x|=n$, and $|p|\le k(n)$.
\item[] \emph{Completeness.} If $U^O(p)$ halts with final output
$x\in\{0,1\}^n$ within $t(n)$ steps, then some permitted $\pi$
satisfies $V^O(1^n,p,x,\pi)=1$.
\item[] \emph{Global list soundness.} For each fixed $(n,p)$,
\[
 \left|\{x\in\{0,1\}^n:\exists \pi,\ |\pi|\le q(n),\
                 V^O(1^n,p,x,\pi)=1\}\right|\le L(n).
\]
This bounds the union over all permitted witnesses, not the output of
one branch or one run.
\end{description}
\end{definition}

\begin{theorem}[List-cover conditional separation]
If this endpoint list cover exists, then $\PP^O\ne\NP^O$.
\end{theorem}
\begin{proof}
Let $E$ be the language whose $n$-bit members are the accepted outputs
over all $p$ of length at most $k(n)$ and permitted witnesses.
Because $k(n)\le n$ and $q$ is polynomial, $E\in\NP^O$.
There are $2^{k(n)+1}-1$ possible programs, so
\[
 |E\cap\{0,1\}^n|
 \le L(n)(2^{k(n)+1}-1)
 \le 2^n-L(n)<2^n.
\]
Thus $\overline E$ has a member at every sufficiently large length.
By completeness, any such member satisfies
\[
 K_U^{t,O}(x)>k(n),\qquad
 K_U^{t,O}(x)\ge n-\ell(n).
\]
After deleting finitely many short lengths, $\overline E$ is an infinite
subset of $R_\ell^{t,O}$. Since $L$ is polynomially bounded,
$\ell(n)=O(\log n)$, so the preceding immunity theorem applies.
If $\PP^O=\NP^O$, then $E$ and its complement are in $\PP^O$,
contradicting that immunity.
\end{proof}

\begin{corollary}[Historical two-valued case]
Taking $L(n)=2$ gives $\ell(n)=1$ and $k(n)=n-2$.
A two-valued polynomially witnessed cover of all these endpoints implies
$\PP^O\ne\NP^O$.
\end{corollary}
\begin{corollary}[Collapsing-oracle obstruction]
For every oracle $O$ with $\PP^O=\NP^O$, no endpoint list cover above
exists. Baker--Gill--Solovay \cite{bgs1975} supplies such oracles.
\end{corollary}
This is a theorem about matched relativizations. It does not establish
nonexistence at the empty oracle. At $O=\varnothing$, existence of the
cover remains an open sufficient condition for $\PP\ne\NP$.

\begin{obligation}[Ordinary bounded-list construction]
Fix a concrete superpolynomial $t$ and construct a verifier meeting all
three conditions for ordinary $U$. Specify the witness polynomial and
prove the global list bound. A succinct computation notation, local
transition checker, or short source program alone is not this construction.
\end{obligation}
The time function is load-bearing. The independent cost-certificate
audit gives an enumeration obstruction when $t$ is sufficiently large
relative to the verifier's witness and verification bounds. No
all-superpolynomial-$t$ construction is asserted here.

\section{Endpoint multiplicity is not accepting-path multiplicity}
For a fixed verifier $V$, define the partial multivalued NP function
\[
 f_V(1^n,p)=\{x:\exists\pi,\ |x|=n,\ |\pi|\le q(n),
                       V(1^n,p,x,\pi)=1\}.
\]
It is realized by an NP transducer that guesses $(x,\pi)$, verifies,
and outputs $x$. Thus global list soundness is exactly a bound on the
number of \emph{distinct output values}. It is not a bound on accepting
paths and does not place the endpoint language in $\mathsf{UP}$ or
$\mathsf{FewP}$.

\begin{lemma}[Witness-union gate]
Let $R\subseteq X\times\Pi$. If every witness fiber
$R_\pi=\{x:(x,\pi)\in R\}$ has size at most $b$, then
\[
 \left|\bigcup_{\pi\in\Pi}R_\pi\right|
 \le \min\{|X|,b|\Pi|\}.
\]
The bound is tight whenever $X$ contains $b|\Pi|$ elements.
Consequently, witnesses of length at most $q$ yield only the bound
$b(2^{q+1}-1)$ absent an overlap theorem.
\end{lemma}
\begin{proof}
The union bound gives the inequality. Equality is obtained by choosing
pairwise disjoint fibers of size $b$.
\end{proof}

\begin{theorem}[Fiber-type quotient gate]
For a finite witness set $\Pi$, let
\(R_\pi=\{x:V(z,x,\pi)=1\}\), put
\[
 b=\max_{\pi\in\Pi}|R_\pi|,
 \qquad
 c=\bigl|\{R_\pi:\pi\in\Pi\}\bigr|.
\]
Then
\[
 \left|\bigcup_{\pi\in\Pi}R_\pi\right|\le bc.
\]
The bound is tight.  Duplicating witnesses or appending ignored witness
bits changes neither $c$ nor the endpoint union.  Thus a polynomial bound
on the number of distinct fiber types, together with a polynomial
per-fiber bound, is a padding-invariant sufficient condition for global
list soundness.
\end{theorem}
\begin{proof}
Choose one representative from each distinct fiber.  Their union is the
original union, and at most $c$ representative sets of size at most $b$
occur.  Pairwise disjoint representatives of size $b$ attain equality.
Duplicating an existing fiber creates no new fiber type and no new endpoint.
\end{proof}

This theorem isolates a useful construction target, not an automatic
compression.  Let the distinct fiber types have multiplicities
$a_1,\ldots,a_c$, let $W=|\Pi|$, and define the exact collision count
\(operatorname{Coll}=\sum_i a_i^2\).  Cauchy--Schwarz gives
\[
 c\ge \frac{W^2}{\operatorname{Coll}}.
\]
This is a lower bound on the number of types, the opposite direction from
the required upper bound.  If a construction additionally proves
$a_i\ge\mu$ for every type, then $c\le W/\mu$ and the endpoint union is at
most $bW/\mu$.  Ignored witness padding can enlarge every $a_i$ without
changing semantics, so multiplicity must be normalized or tied to a
construction-specific witness convention.

\begin{lemma}[Truth-preserving selector gate]
Let $T_z\subseteq F_z$ be the required true endpoint set, with
$|T_z|\le1$, and let $F_z$ be the set accepted by a permissive verifier.
Any selector that sees only $F_z$ and returns a strict subset of some
$F_z$ fails to preserve $T_z$ for one admissible singleton
$T_z\subseteq F_z$.
\end{lemma}
\begin{proof}
If $C(F_z)\subsetneq F_z$, choose
$x\in F_z\setminus C(F_z)$ and take $T_z=\{x\}$.  Then
$T_z\nsubseteq C(F_z)$.
\end{proof}
Consequently a generic refinement of the accepted NPMV endpoint function
does not by itself retain the actual universal-machine output.  The
canonicalizer needs a source/program binding predicate or another theorem
that distinguishes the required endpoint.  Hemaspaandra--Naik--Ogihara--
Selman show that generic unique refinements of NPMV functions have
polynomial-hierarchy collapse consequences \cite{hnos1996}.  Inside a
$P=NP$ contradiction argument that qualitative collapse is already
available, but it supplies neither this truth-preserving binding nor the
needed simultaneous exponent bounds.

Already with two outputs and two witnesses, the fibers
$R_{\pi_0}=\{x_0\}$ and $R_{\pi_1}=\{x_1\}$ are individually
single-valued while their existential union has two values. Hence a
per-certificate singleton claim cannot replace global list soundness.

\begin{lemma}[Witness padding separates the notions]
Appending $d$ ignored witness bits preserves $f_V(1^n,p)$ exactly and
multiplies every accepting-witness count by $2^d$.
\end{lemma}
\begin{proof}
Accept precisely when the original witness prefix is accepted, ignoring
the new suffix. The endpoint image is unchanged, while every old
witness has all $2^d$ suffix extensions.
\end{proof}
Therefore output multiplicity is invariant under harmless witness
padding whereas $\mathsf{FewP}$-style path multiplicity is not.

Lexicographic canonicalization does not repair this for free. The
predicate that an accepted witness has no smaller accepted witness has
a coNP universal subcondition for a polynomial-time verifier. Under
the contradiction hypothesis $\PP=\NP$ it becomes polynomial-time
decidable, but with an uncontrolled exponent; it is neither an
unconditional NP-verifier construction nor a charged endpoint theorem.
Generic reductions from multivalued to uniquely valued NP functions are
known to have collapse consequences \cite{hnos1996}; they cannot be
silently imported as a routine normalization.

\section{Why sparse-set theorems do not follow from a per-program list}
Mahaney's theorem \cite{mahaney1982} concerns a globally sparse
NP-hard language under polynomial-time many-one reductions. The
PFCERT list condition instead bounds, for each fixed $(n,p)$, the number
of endpoint values. These are different density statements.

\begin{lemma}[Fiber sparsity need not be joint sparsity]
There is a polynomial-time relation with exactly one endpoint for every
program but exponentially many accepted program--endpoint pairs at
infinitely many total lengths.
\end{lemma}
\begin{proof}
For each $(n-1)$-bit string $p$, accept only $x=0p$. Each program fiber
has size one, yet there are $2^{n-1}$ accepted pairs $(p,x)$ of total
length $2n-1$ (up to the fixed encoding overhead).
\end{proof}
Thus neither Mahaney's theorem nor an immunity theorem applies to the
joint graph without a new reduction and global-density theorem.
Likewise, current range-avoidance results under generator assumptions
\cite{renwangzhong2026} and exact search-to-decision results in declared
regimes \cite{hkloliveira2024} do not construct the ordinary all-input
bounded-list cover. All arguments in this section relativize; they
remove invalid transfers but do not supply the missing nonrelativizing
hinge.


\section{A linear-sharp source-probe guard}
Consider the proper 3-CNF switch family in which forcing gadgets set
\(x_1,\ldots,x_r,y,z\) to zero, the satisfiable member contains
\[
 C_i^0=(\neg x_i\vee y\vee z)
\]
in every switch slot, and member \(F_j\) replaces only slot \(j\) by
\(C_j^1=(x_j\vee y\vee z)\).  The satisfiable member and every
unsatisfiable one-slot flip have opposite SAT values.

\begin{theorem}[Source-probe tightness and matched-oracle guard]
Every deterministic self-contained exact transporter must query all
\(r\) switch slots on the satisfiable member.  This bound is tight:
a direct scan decides the family and emits one Boolean value in
\(O(r)\) time and \(O(\log r)\) work bits.  Both claims hold relative
to any common oracle \(O\).
\end{theorem}
\begin{proof}
If slot \(j\) is unqueried, the deterministic transcript and output are
identical on the satisfiable member and its one-slot unsatisfiable flip.
Conversely, scan the signs and accept exactly when all are negative.
The proof neither examines nor restricts oracle answers, so attaching
the same oracle everywhere changes nothing.
\end{proof}

This theorem classifies the clause-probe argument as an exact near-linear
input-access guard, not as the missing all-algorithm lower bound.  It is
achieved by an easy algorithm on the same family and relativizes.
Under the hypothesis \(\SAT\in\PP\), decide-then-emit also gives a
constant-length semantic successor while retaining the unknown solver's
time and space.  Therefore strict output shrinkage, history
distinguishability, or source grounding alone cannot be promoted to a
favorable charged successor or a nonrelativizing hinge.

\section{Compatibility-rank guard inside the exact clean core}
Let $B$ be the paired-incidence formula of $K_{3,3}$: one Boolean variable
for each edge and, at each vertex, the positive and negative clause on its
three incident edges.  Then $B$ is a proper unique Monotone
$3$-Sat-$(2,2)$ formula.  A perfect matching satisfies it.

For the row-major edge indexing and split
\[
 X=\{2,4,6,7,8\},\qquad Y=\{0,1,3,5\},
\]
the semantic compatibility matrix $H_B(a,b)=B(a,b)$ has shape
$32\times16$ and full column rank 16 over both $\mathbb F_2$ and
$\mathbb R$.  If $F_k$ is the conjunction of $k$ variable-disjoint
canonical copies, then
\[
 H_{F_k}=H_B^{\otimes k},\qquad
 \operatorname{rank}H_{F_k}=16^k=2^{4k}.
\]
The indexed encoding has length $N_k=\Theta(k\log k)$, so the rank is
$2^{\Theta(N_k/\log N_k)}$.

This does not lower-bound the SAT decision cost.  A uniform decider compares
the source with the canonical block template in $O(N_k)$ time and
$O(\log N_k)$ writable workspace and returns SAT; a fixed matching is a
witness for every block.  Off the family it may invoke any total exact
fallback.  Under a hypothetical polynomial SAT decider the same
recognize-first construction remains polynomial on all inputs.  Therefore
compatibility rank, assignment-split communication, matrix materialization,
and tensor state are not forced resources of a one-bit SAT execution.
The unrestricted worst-case target, an arbitrary-decider compiler, and a
nonrelativizing/nonalgebrizing hinge remain open.  P versus NP remains open.

\section{Clause-local versus incidence-local reduction access}
For a length-$L$ raw occurrence stream in which a designated key occurs
exactly three times and one occurrence is at position $p$, let the query
output its sign-specific occurrence rank.  In the deterministic stateless
field-probe model the exact worst-case query complexity is
\[
 \min\{p,L-p-1\},
\]
apart from the constant-output boundary cases.  The upper bound scans the
smaller side.  For the lower bound, answer every query by ``not the key'';
before that many probes at least one unqueried position remains on each side
and one side has two, permitting two consistent completions with different
ranks.

This separates clause locality from incidence locality, but not P from NP.
On a Monotone 3-Sat-$(3,3)$ source, two native tag bits per literal label the
three occurrences of each variable and sign.  The tag lift adds only $12n$
bits for $n$ variables, hence preserves indexed encoding length within a
constant factor, and the corrected Darmann--D\"ocker clone is then selected
by constant local arithmetic.  Preprocessed succinct rank/select offers a
second model-distinct bypass once its storage and construction are charged.

Thus the raw scan theorem is a black-box composition guard, not an
encoding-invariant or arbitrary-decider lower bound.  The remaining first
leg must provide both local clause access and native/charged incidence tags.
That bi-local reduction, the arbitrary-exponent bootstrap, UCORE-TIME1, and
P versus NP remain open.

\section{Annotation-semantics trilemma}
For any language \(L\), annotation map \(\tau\), and locally projectable
encoding \(\langle x,a\rangle\), distinguish
\[
 L_{\rm erase}^{\tau}=\{\langle x,a\rangle:x\in L\},
 \qquad
 \Pi_L^\tau=\{\langle x,\tau(x)\rangle\},
 \qquad
 L_{\rm valid}^{\tau}=\{\langle x,a\rangle:x\in L\wedge a=\tau(x)\}.
\]
A decider for \(L\) gives a uniform decider for the erased total language by
simulating on the virtual \(x\)-projection and never reading \(a\).  The same
machine solves the correctly annotated promise problem without checking the
promise.  The validated object is a different total language: on
\(x\in L,a\ne\tau(x)\), erasure accepts and validation rejects.  Therefore a
validation lower bound transfers to ordinary SAT only through an explicit
compiler producing valid annotations, with construction, local access,
storage and workspace charged in the same realization.

This is \texttt{ANNOTATION-SEMANTICS-TRILEMMA1} at its declared interface.
Unqualified \texttt{TAG-VALIDATION-TO-UCORE1} is
\textbf{FALSIFIED/RETIRED}.  Occurrence tags remain useful external reduction
infrastructure; their validation is not forced on a one-bit SAT decider.
\texttt{JMV-NATIVE-TAG1}, local tagged Monotone \((3,3)\) normalization,
SAME\_REALIZATION endpoint composition, UCORE-TIME1, and P versus NP remain
open.

\section{Exact black-box query-scale dichotomy}
Let a uniform deterministic extensional transducer make at most
\(N^{\eta+o(1)}\) exact SAT queries, use \(N^{r+o(1)}\) charged local
time and \(N^{p+o(1)}\) local workspace.  Let \(q_E(N)\) be the maximum
length of an actually reached query over length-\(N\) inputs and set
\(\sigma_E(N)=\log q_E(N)/\log N\).

Assume SAT is in P.  If \(\sigma_E(N)\to0\), fixing any polynomial exact
SAT solver with time exponent \(K\) evaluates each query in
\[
  q_E(N)^{K+o(1)}=N^{o(1)}
\]
time and space, so the full ordinary source computation has time
\(N^{\max\{r,\eta\}+o(1)}\) and workspace
\(N^{p+o(1)}+N^{o(1)}\).  If \(\sigma_E(N)\not\to0\), some fixed
positive query scale occurs infinitely often; an extensionally identical
polynomial delay wrapper then defeats every exponent quantified before all
polynomial SAT implementations.

This is \texttt{BLACKBOX-QUERY-SCALE-DICHOTOMY1}.  The vanishing-scale half
is the zero-query-exponent case of the existing direct-wrapper simulation;
the positive-scale half is MODELBRIDGE-0015.  Together they are exhaustive.
Thus \texttt{SUBPOLY-BLACKBOX-IS-INTERMEDIATE1} is
\textbf{FALSIFIED/RETIRED}: under \(P=NP\), polynomially many
\(N^{o(1)}\)-length exact SAT calls yield an ordinary source algorithm, not
a third oracle-relative hardness regime.  Constructing such an algorithm in
an FLMV-forbidden time--space region remains \textbf{OPEN} and
separation-strength.  P versus NP remains open.


\section{All-branch restriction fanout at encoded length}
Let \(F\) be a duplicate-free proper exact-\((2,2)\) 3-CNF with
\(m=4n/3\) and canonical indexed length \(N=\Theta(n\log n)\). For a block
\(S\), explicitly emit all \(2^{|S|}\) residuals. At least \(m-4|S|\)
clauses avoid \(S\) and survive unchanged in every residual. Since residual
degree is at most four, \(r\) such clauses use at least \(3r/4\) live
variables and require \(\Omega(r\log(r+2))\) bits even after compact
renaming.

Therefore, if every child has length at most \(N^h\) for fixed \(0<h<1\),
then
\[
 |S|\ge n/3-o(n),\qquad
 2^{|S|}=2^{\Omega(N/\log N)}.
\]
This is \texttt{ALL-BRANCH-ENCODED-FANOUT1}, PROVED/EXACT for the declared
explicit pass. It retires \texttt{ALL-BRANCH-SOURCE-PASS1} only as a
polynomial-cost strict power contraction. A pruned source-dependent pass,
implicit representation with charged access, and arbitrary one-bit SAT
computation remain open. Adaptive self-reduction uses current-scale calls
and retains the unknown solver exponent. Compatibility size is not work:
the affine control has all sixteen block assignments compatible but a
four-row Gaussian representation. P versus NP remains open.

\paragraph{Verification.}
The FLMV rational inequality and bounded-list counts were checked by exact
integer/rational arithmetic on small parameters as a sanity check.  The
fiber-type quotient, collision direction and selector-binding controls were
also exhaustively checked for all four-by-four relations; the proofs above
do not rely on those finite checks. An independent REDTEAM
agent reviewed the oracle-model and witness-bound guards.
No full-repository build or proof-assistant check was performed.


\section{Pruned strict-power restrictions carry linear partial witnesses}
Let $F$ be a duplicate-free proper exact-$(2,2)$ 3-CNF on $n$ variables with canonical indexed length $N=\Theta(n\log n)$.  If a literal restriction $\rho$ leaves a residual of length at most $N^h$ for fixed $0<h<1$, then only $o(n)$ clauses may survive.  Since assigning one variable can delete at most four source clauses and $F$ has $4n/3$ clauses,
\[
 |\mathrm{dom}(\rho)|\ge n/3-o(n).
\]
If $F$ is satisfiable and $F\upharpoonright\rho$ remains satisfiable, $\rho$ is an extendable partial witness of that size.  Explicitly serializing it costs $\Omega(n)=\Omega(N/\log N)$ bits.  This is \texttt{PRUNED-CHILD-LINEAR-ASSIGNMENT1} and \texttt{PRUNED-YES-PARTIAL-WITNESS1}, PROVED/EXACT at the declared explicit-restriction interface pending independent review.  It does not force an arbitrary one-bit decider to materialize the restriction: \texttt{PRUNED-WITNESS-FORCES-DECIDER-WORK1} is FALSIFIED/RETIRED, while \texttt{IMPLICIT-PREDECISION-ACCESS1} remains OPEN.  P versus NP remains open.

\bibliographystyle{plain}
\bibliography{lower_bound_refs}
\end{document}


